\subsection{A normal form for every nondegenerate binary residual}
Let $E$ have binary basis matrix $V$ and nonsingular Gram matrix
$G=V^tV$. In dual Fourier coordinates put
$q(z)=\operatorname{wt}(VG^{-1}z)\pmod4$.
Its polarization is $2z^tG^{-1}w$. The same argument applies to mixed
signs on orthogonal summands, since reversing a sign preserves this
polarization. Write $S_q=\sum_z i^{q(z)}$. Completing the square gives
\[
\begin{aligned}
K(x,y)&=2^{-r}\sum_z i^{q(z)}(-1)^{z^t(x+y)}\\
 &=2^{-r}S_q\,i^{-q(Gx)}(-1)^{x^tGy}i^{-q(Gy)}.
\end{aligned}
\]
The scaled Walsh kernel in the middle is a binary permutation of
$H_0^{\otimes r}$, where
\[
 H_0=\tfrac12\begin{pmatrix}1&1\\1&-1\end{pmatrix}
     =bSCS,\qquad b=(1-i)/2,\quad S=\operatorname{diag}(1,i).
\]
The remaining scalar $S_qb^r$ is a fourth root of unity. To verify
this, decompose the nondegenerate quadratic module into odd
one-dimensional and alternating two-dimensional blocks. Their Gauss
sums are $1\pm i$ and $\pm2$, respectively; multiplying by $b$ and
$b^2$ gives a fourth root in each case. This is a proof device, not
an orthonormal-basis assumption on $E$.

Consequently one rank-$r$ child suffices, with binary address changes,
diagonal fourth-root phases, and one scalar fourth-root phase.
On $f$ parallel columns it is one $C^{\otimes rf}$ child, the scalar
raised to power $f$, and columnwise quadratic phases modulo four.
The retained packed address operations evaluate these fixed formulas.
These wrappers have modulus one and introduce no denominator loss.

\subsection{The selected Gaussian-dyadic producer}
Use paired triple exclusion with base threshold two, complete unsplit
disjoint and pair-star sums, and sharing by actual ordinary support.
For an ordinary support with at least two common points use the span
of its triple indicators; these form an orthonormal common-pair star.
Otherwise use the coordinate space of its union. Support inclusion
preserves these labels: losing the common pair enters a containing
coordinate space. Source lines use ordinary triple containment.

Retain every full pair-star $P_{uv}=\sum_{w\ne u,v}X_{uvw}$ and form
$B_u=\sum_{v\ne u}P_{uv}=2\sum_{T\ni u}X_T$.
For even $h$ label these central additions by the nondegenerate
hyperplane $E_u=\{x:\sum_{j\ne u}x_j=0\}$; every contributing star
lies in it. Retain the global total $T$ with full-space label.
The central correction at target $S$ is
\[
 \frac14\sum_{u\in S}B_u-\frac12T
       =\sum_U\frac{|S\cap U|-1}{2}X_U.
\]
Adding half the disjoint sum and subtracting half the intersection-two
sum gives exactly $X_S$. All coefficients are Gaussian dyadic.
The transparent retained-output wrapper and oriented carrier compiler
apply to these nested nondegenerate labels. The $h$ hyperplane centers
lose $h(h-1)$ and the global center loses $h$, so loss is $h^2$.
The selected counts, including all point centers, are
\[
\begin{array}{c|rrrrrr}
h&v_h&c_h&q_h&|\mathcal M_h|&R_h&\text{loss}\\\hline
32&4960&122512&19873&19241&123144&1024\\
30&4060&99153&16271&15873&99551&900\\
40&9880&252023&39561&37642&253942&1600
\end{array}
\]
Their internal histograms satisfy $\sum_r rH_{h,r}=hR_h+2h^2$.
They are distinct from the bit histograms.

\subsection{Endpoint normalization and the complex characteristic}
For nondegenerate $U$, Fourier phases give
\[
 C_U^2=X_{P_U\mathbf1},\qquad
 C_{U^\perp}C_U^{-1}=C_FX_{P_U\mathbf1}.
\]
Put $v=P_U\mathbf1$. On coordinates in $v$, this endpoint has
$C^{-1}$ instead of $C$. The identity $C^{-1}=-iZCZ$ shows that
conjugating by $Z_v$ and multiplying by $i^{\operatorname{wt}(v)}$
normalizes it to $C_F$; for $f$ columns the exponent is
$f\operatorname{wt}(v)$. These are diagonal wrappers without a child.

Choose stage-two source $C_{D_0}$ and sink $C_{D_0^\perp}$, where
$D_1=D_0\mathbin{\perp} A$ and $\dim A=b$. The same refined copy
schedule removes the entrance on every auxiliary role. With
$R=D_1^\perp$, the exit is $C_RC_{D_0}^{-1}$ on
$E=R\mathbin{\perp}D_0=A^\perp$. Its rank is $m-b$ and its
nondegenerate polar form is unchanged by the mixed signs. The normal
form above executes it as one child. Arbitrary auxiliary restoration
justifies the endpoint normalization on the entire $B_2$ bank.

Use the complex counts to define $B_i=NR_{h_i}/v_{h_i}$, and now
\[
\begin{aligned}
W&=2N+B_2+B_3=10390178240000,\\
L&=\sum_i(N/v_{h_i})h_i^2=117399987200,\\
s&=Wm-2N+2L=398982681296998400.
\end{aligned}
\]
The whole-residual macro list is
\[
\begin{array}{c|rrrrr}
t&38328&38360&38370&37401&899\\\hline
n_t&B_1&B_3-B_1&B_2&2N&2N.
\end{array}
\]
Add $(N/v_h)H_{h,r}$ rank-$r$ children for each positive internal rank.
Leave the remaining physical-data rank, 39393978624000, as singleton
calls. The resulting list has weighted rank exactly $s$.
For $a_c=3/125000$ and $\sigma=1-a_c$, the rational logarithm bounds
above and $e^u\le(1-u)^{-1}$, $0\le u<1$, give
\[
 \Psi_c(\sigma)\le\sum_t\frac{n_tt}{Wm(1-a_c\ell_t)}
 <1-\frac{792981575160764}{10^{23}}.
\]

\subsection{The dependency-path guard}
Before gauging, nested producer labels telescope along every arithmetic
path, including wire switches at common-frame scalar gates. Retained
center drops occur on one return frontier per stage, crossed at most
once by a path, and each costs at most that axis dimension. Carrier
reuse adds only increasing continuations. This gives
$q_0=m+2(a+b+k)$. Gauging removes an early edge and adds $ab-b$ at
a late stage-two edge. A path meets at most one such edge, since it
terminates that stage-two role, which is not shared across stages.
Thus the conservative total rank budget is
\[
 q=m+2(a+b+k)+ab=39564.
\]
Every child except the four large macros is at most 899; two large
macros cannot fit within $q$. Set $\rho=19/10$, $\gamma=9997/10000$.
The integer comparisons $m^{19}>500000000^{10}$ and
$899^9<460^{10}$ give $m^\rho>500000000$ and
$899^{\rho-1}<460$.
If a path has one large child $t$, its moment is at most
\[
 (t/m)^\rho+(q-t)899^{\rho-1}/m^\rho.
\]
This expression increases over the large-child range, since its
numerator derivative is $\rho t^{\rho-1}-899^{\rho-1}>0$.
Use $t=m-b$ and the concavity bound
$(1-x)^{19/10}\le1-19x/10+9x^2/10$ to obtain
\[
 \sum_{\rm path}(r/m)^\rho
 <1-\frac{19}{12800}+\frac9{10\cdot1280^2}
       +\frac{1194\cdot460}{500000000}
 =\frac{51180270301}{51200000000}<\gamma.
\]
Without a large macro the bound is $227493/6250000<\gamma$.

For coefficient-depth accounting take $E_0=128(W+m+1)^3$.
There are at most $s$ positive-rank residual calls and $W$ endpoint
normalizations. Combining the quadratic phases with the $S$ factors,
each uses at most two diagonal phase passes and one scalar phase.
The inherited gate and remainder count is therefore bounded conservatively by
\[
 36W^3+16s+16W+8m+4<E_0.
\]
Binary address changes only permute coefficients; descriptor evaluation
is paid by the retained packed record work and does not increase
coefficient depth.
Stop below width $d^\beta$, where $\beta=1/1000$.
For $e=mf+r$, $0\le r<m$, recurse on $tf$ and treat the remainder
individually. For sufficiently large $d$, a leaf descending from an
internal node has width at least $d^\beta/(2m)$. Choose
$C_{\rm dep}\ge16m$ and
$(1-\gamma)C_{\rm dep}\ge E_0+16m+1$.
The retained elementary-kernel bound $A(e)\le8e$ and strong induction
using the path moment imply
\[
 A(e)\le C_{\rm dep}d^{-\beta(\rho-1)}e^\rho.
\]
At a leaf the lower width bound proves this inequality, while at an
internal node the unused $1-\gamma$ fraction covers local work.
Small roots are handled directly. With $\zeta=1/10000$ take
\[
 C_1=\rho-(\rho-1)\beta+\zeta=\frac{1187}{625},\quad
 C_0=\left\lceil128m(1+1/\zeta)C_{\rm dep}\right\rceil,
 \quad\Delta=\lceil C_0d^{C_1}\rceil.
\]
The retained exact denominator/magnitude-depth induction gives
$p+\Delta$ fractional bits and $\Delta+3$ integer/sign bits.
All selected scalar coefficients are Gaussian dyadic and the new
wrappers have modulus one, so that induction still applies.
For $\epsilon C_1<1$ their total width is $O(p)$.
\subsection{Rows, stopping, and transfer to the layer interface}
Let $t_*<m$ be the largest child and compute
$D(e)=0$ for $e<m$ and $D(e)=1+D(t_*\lfloor e/m\rfloor)$ otherwise.
It is $O(\log e)$ and bounds every descendant depth.
With the retained compact-control parameters, reserve
\[
 q_0=\lceil\log_2W\rceil D(d),\qquad Q_0=q_0+q_F+q_B.
\]
Process the first $q_0+q_F$ and last $q_B$ selected axes individually;
if at most $Q_0$ axes exist, process them all individually. Otherwise
the middle active interval is consecutive. Only the $q_0$ row chunks
supply row indices, separate from the two compact fields. Pad their
count once to a multiple of $W^{D(d)}$, by a factor below two.
Every child retains complete compact and spectator ranges and has
exactly $1/W$ of its parent's volume. No internal padding is needed.
The reserved-axis bound remains
\[
 Q_0=O(\log(2d)+d^{\max\{1-c,0\}}\log p+1).
\]
Remainders commute with the completed tensor map and cost $O(V)$ per
node. The fixed-tape depth-first scheduler and complete-row argument
of the bit construction apply with the retained compact descriptors.

Since $\Psi_c(\sigma)<1$, the total leaf $\sigma$-weighted width,
including the volume factor $W^{-j}$ at depth $j$, is at most
$d^\sigma$. Thus leaf work is
$O(d^{\sigma+\beta(1-\sigma)})$. Internal compact operations cost
$O(G^\tau e^\tau+1)$ per current volume. Summing the strict moments
gives the retained complete normalized layer bound
\[
 O\!\left((\log p)^3\left[
 d^\chi+d^{\sigma+\beta(1-\sigma)}+d^{\max\{1-c,0\}}+1
 \right]\right),\qquad
 \chi=\tau+(1-\beta)\max\{\sigma-\tau,0\}.
\]
For $\sigma\le\tau$, use $\Psi_c(\tau)\le\Psi_c(\sigma)$;
otherwise use $e^\tau\le d^{\beta(\tau-\sigma)}e^\sigma$ at each
internal node. These arguments do not assume equal leaf depth.
The usual strict conditions on $\lambda,\lambda'$ therefore give
the retained $O(Vd^{\lambda'})$ layer interface, now with the new
moment and the new $C_{\rm dep}$.
